\documentclass[a4paper]{article}

% Español (México) con XeLaTeX: fontspec para fuentes Unicode, polyglossia
% para la división silábica y los nombres del documento en la variante mexicana.
\usepackage{fontspec}
\usepackage{polyglossia}
\setdefaultlanguage[variant=mexican]{spanish}
% Los nombres chinos van como «pinyin (original)»: xeCJK da a los caracteres
% Han su propia fuente; sin ella XeLaTeX los omite con un aviso.
\usepackage{xeCJK}
\setCJKmainfont{FandolSong-Regular.otf}
% polyglossia no traduce los nombres de listings.
\AtBeginDocument{\providecommand{\lstlistingname}{}\renewcommand{\lstlistingname}{Listado}%
  \providecommand{\lstlistlistingname}{}\renewcommand{\lstlistlistingname}{Índice de listados}}
\usepackage{amsmath, amssymb}
\usepackage{graphicx}
\usepackage[margin=2.5cm]{geometry}
\usepackage[most]{tcolorbox}
\usepackage{etoolbox}
\usepackage{listings}
\usepackage{hyperref}
\usepackage{booktabs}
\usepackage{array}
\usepackage{tabularx}
\usepackage{longtable}
\usepackage{float}
\usepackage{tikz}
\usetikzlibrary{arrows.meta,positioning,calc}

\newcolumntype{Y}{>{\raggedright\arraybackslash}X}

\newtcolorbox{knowledgebox}[1]{
    enhanced, colback=blue!5!white, colframe=blue!75!black, colbacktitle=blue!75!black,
    coltitle=white, fonttitle=\bfseries, title={#1}, attach boxed title to top left={yshift=-2mm, xshift=2mm},
    boxrule=1pt, sharp corners
}

\newtcolorbox{importantbox}[1]{
    enhanced, colback=yellow!10!white, colframe=yellow!80!black, colbacktitle=yellow!80!black,
    coltitle=black, fonttitle=\bfseries, title={#1}, sharp corners
}

\newtcolorbox{warningbox}[1]{
    enhanced, colback=red!5!white, colframe=red!75!black, colbacktitle=red!75!black,
    coltitle=white, fonttitle=\bfseries, title={#1}, sharp corners
}

\lstset{
    language=Python,
    basicstyle=\ttfamily\small,
    keywordstyle=\color{blue},
    stringstyle=\color{red!60!black},
    commentstyle=\color{green!60!black},
    breaklines=true,
    frame=single,
    numbers=left,
    numberstyle=\tiny\color{gray},
    captionpos=b,
    extendedchars=true
}

\newcommand{\notetitle}{Notas de entrevista: Zhao Chenyang (赵晨阳) habla de SGLang, la ola de abandono escolar, AI Coding y la comunidad de código abierto}
\newcommand{\noteauthors}{Elaboradas a partir de Ungrounded 不着边际 EP02}
\newcommand{\notedate}{2026-05-08}
\newcommand{\videochannel}{Ungrounded 不着边际}
\newcommand{\videopublishdate}{2026-05-06}
\newcommand{\videoduration}{01:07:50}
\newcommand{\videourl}{https://www.bilibili.com/video/BV1oRRyBREPU/}
\newcommand{\videocoverpath}{cover.jpg}
\newcommand{\repourl}{https://github.com/hqhq1025/ai-course-notes}

\usepackage{fancyhdr}
\pagestyle{fancy}
\fancyhf{}
\fancyfoot[C]{\thepage}
\fancyfoot[R]{\footnotesize\color{gray}\href{\repourl}{AI Course Notes}}
\renewcommand{\headrulewidth}{0pt}
\renewcommand{\footrulewidth}{0pt}

\begin{document}

\begin{titlepage}
\centering
{\Large Notas de entrevista\par}
\vspace{1.2cm}
{\huge\bfseries \notetitle\par}
\vspace{0.8cm}
{\large \noteauthors\par}
\vspace{0.3cm}
{\large \notedate\par}
\vspace{1.2cm}

\ifdefempty{\videocoverpath}{
{\small Al generar las notas, indique la ruta local de la portada del video.\par}
}{
\includegraphics[width=0.82\textwidth,height=0.45\textheight,keepaspectratio]{\videocoverpath}\par
}

\vfill
\begin{tcolorbox}[width=0.9\textwidth, colback=black!2!white, colframe=black!60, sharp corners]
\textbf{Autor o canal del video}: \videochannel\par
\textbf{Fecha de publicación}: \videopublishdate\par
\textbf{Duración del video}: \videoduration\par
\textbf{Enlace del video}: \href{\videourl}{\nolinkurl{\videourl}}\par
\textbf{Página del proyecto}: \href{\repourl}{\nolinkurl{\repourl}}\par
\end{tcolorbox}
\end{titlepage}

\tableofcontents
\newpage

\section*{Introducción}

Esta conversación parece, en la superficie, tratar sobre Zhao Chenyang (赵晨阳): licenciatura en Tsinghua, quince meses de doctorado en UCLA, participación en SGLang, y el paso de un proyecto de código abierto a fundar una empresa de motores de inferencia; cuando se publicó el video, RadixArk acababa de anunciar una ronda semilla de cien millones de dólares. Si uno se queda solo en los titulares, es fácil entenderla como «la historia de un joven más en la ola de abandonos universitarios de Silicon Valley». Pero lo que realmente vale la pena conservar de toda la charla no es el abandono en sí, sino la pregunta más profunda a la que apuntan juntos el abandono, SGLang, AI Coding y la comunidad de código abierto: una vez que la IA reduce con rapidez el costo de ejecución, ¿qué responsabilidades le quedan a la persona que no puede delegar?

Zhao Chenyang insiste una y otra vez en que «hay que hacer las cosas de manera esencial». Esta idea recorre toda la charla. Cuando habla del mundo académico, critica la inercia de tener un martillo y salir a buscar clavos, de formular una evaluation en el vacío; cuando habla de dejar el doctorado, le preocupa si la applied research realmente necesita someterse a la validación del mercado y de la ingeniería; cuando habla de AI Coding, subraya que, una vez que producir código se vuelve barato, las personas necesitan hacerse cargo del design, del review, de decir que no (say no) y de asumir las consecuencias; cuando habla de la comunidad de código abierto, lo que en realidad le importa no es «cuántas contribuciones» se hacen, sino cómo un proyecto mantiene sus límites de responsabilidad, su sentido estético y su orden de colaboración una vez que el código generado por IA se vuelve abundante.

Por eso, estas notas no siguen el orden cronológico, sino que se reorganizan en torno a seis temas: la posición ecológica de los motores de inferencia, el cambio de paradigma de investigación detrás de la ola de abandonos, el research taste en la era de la IA, el impacto de AI Coding sobre la gobernanza del código, el riesgo cognitivo que trae la IA como amplificador, y cómo la comunidad de código abierto redefine la contribución en la era del código generado.

\begin{importantbox}{Tres hilos conductores para leer estas notas}
El primer hilo es la «esencialidad»: la applied research necesita una application, y la fundamental research necesita un sistema de problemas coherente consigo mismo; no basta con optimizar contra un blanco ficticio. El segundo hilo es el «límite de responsabilidad»: la IA abarata la ejecución, pero también facilita que se multipliquen los pull requests, el código y el contenido hechos sin pensar. El tercer hilo es la «comunidad y la organización»: un proyecto de código abierto no es solo un repositorio de código, sino que debe registrar contribuciones, filtrar responsabilidades, transmitir información y organizar a las personas para que trabajen juntas.
\end{importantbox}

El video original no tiene subtítulos oficiales; esta transcripción se elaboró a partir de una transcripción local con Whisper y una revisión párrafo por párrafo. El cuerpo del texto no reproduce la transcripción palabra por palabra, sino que busca extraer, en la medida de lo posible, las cadenas de juicio y los puntos de vista transferibles de la entrevista.
\section{La posición de SGLang: la capa intermedia en la arquitectura de cinco capas de la IA}

Zhao Chenyang (赵晨阳) comienza presentándose: licenciatura en ciencias de la computación en Tsinghua, egresado en 2024, cursó quince meses de doctorado en UCLA y después, alrededor de noviembre de 2025, comenzó a emprender, uniéndose a una empresa orientada a motores de inferencia. El presentador dirige enseguida la pregunta hacia SGLang: cómo pasó de ser una organización de código abierto a convertirse en empresa, y qué experiencia tuvo Zhao Chenyang en ese proceso.

Primero explica qué es un motor de inferencia. Cuando un usuario común entra en contacto con la IA, lo que ve son productos de frontend como ChatGPT, Claude, AI PPT, Sora o AI Coding; pero una vez que la solicitud del usuario llega al backend, cómo el servidor programa el modelo, cómo administra el KV cache, cómo cuantiza, cómo maneja distintas formas de modelo, cómo aumenta el rendimiento y reduce la latencia: eso es lo que debe resolver el motor de inferencia. El \texttt{model.generate} de Hugging Face ayuda a entender el proceso de generación más elemental, pero atender solicitudes reales en producción exige una gran cantidad de optimización a nivel de sistema.

Zhao Chenyang recurre a una arquitectura de cinco capas para ubicar la inferencia: en la capa superior están los agentes y las aplicaciones; debajo, los modelos concretos, como Claude, GPT o modelos de código abierto; más abajo, los motores de inferencia, como OpenAI internal engine, TensorRT-LLM o SGLang; más abajo, el hardware, y más abajo aún, infraestructura como los semiconductores y la energía eléctrica. Cuanto más arriba, más rápido cambia; cuanto más abajo, más lento es el ciclo; el motor de inferencia se ubica justo en medio, como la capa de enlace clave entre los fabricantes de modelos y los fabricantes de hardware.

\begin{knowledgebox}{Por qué importa el motor de inferencia}
\begin{tabularx}{\textwidth}{>{\raggedright\arraybackslash}p{0.24\textwidth} Y}
\toprule
Nivel & Problema que debe resolver el motor de inferencia \\
\midrule
Programación de solicitudes & Cómo poner en cola de manera eficiente a distintos usuarios, distintas longitudes de contexto y distintos lotes. \\
KV Cache & En contextos largos y conversaciones de varios turnos, cómo reutilizar, trasladar y liberar la caché. \\
Forma del modelo & Procesos de inferencia distintos, como autorregresivo, difusión o flow matching, tienen necesidades de sistema distintas. \\
Estructura de costos & La latencia, el rendimiento, la memoria de video y el aprovechamiento del hardware determinan en conjunto si el producto puede escalar. \\
Adaptación del ecosistema & Los modelos de la capa superior cambian rápido y el hardware de la capa inferior cambia lento; la capa intermedia debe absorber continuamente los cambios de ambos lados. \\
\bottomrule
\end{tabularx}
\end{knowledgebox}

Lo más valioso que hay que conservar de este pasaje no son los detalles técnicos concretos de SGLang, sino el juicio de Zhao Chenyang sobre la «capa intermedia». En la era de la IA, cuanto más arriba está la capa, más ruidosa es y más fácil que la llene la narrativa del producto; pero lo que en realidad determina si una aplicación puede escalar suele ser el costo, el rendimiento y la confiabilidad de ingeniería de la capa intermedia. El motor de inferencia no genera noticias con la facilidad con que lo hace el lanzamiento de un modelo, ni enfrenta directamente al usuario como una aplicación, pero se convierte en un cimiento importante para que funcione toda la economía de la IA.

\subsection{Resumen de la sección}

La capa de motores de inferencia en la que se ubica SGLang es la capa intermedia entre las aplicaciones de IA, los modelos y el hardware. No responde directamente «qué tan inteligente es el modelo», sino «cómo mantener el servicio del modelo funcionando de manera estable con un costo asumible». La experiencia de Zhao Chenyang es adecuada para discutir la ola de abandono del doctorado precisamente porque no participó en un proyecto académico abstracto, sino en un sistema de código abierto situado en el centro de la estructura de costos de la industria.
\section{La ola de abandono de estudios: no es antiacademia, sino que la applied research se ve obligada a enfrentar el mercado}

El abandono de estudios es el tema más visible de este episodio, pero ni Zhao Chenyang (赵晨阳) ni el conductor lo presentan como una historia de rebeldía juvenil. Zhao Chenyang dice que cursó quince meses de doctorado, y que esos quince meses coinciden en gran medida con el tiempo en que desarrolló SGLang. Después consideró que el timing era correcto, que la industria ofrecía más oportunidades que la academia, y por eso optó por irse. La palabra clave detrás de esta decisión no es «abandonar los estudios», sino «qué tipo de research debe ocurrir en qué lugar».

Él divide la research en applied research y fundamental research. La applied research, precisamente por llamarse applied, debe tener una application: debe someterse de verdad a la prueba del mercado, los usuarios, los sistemas de ingeniería y el valor comercial. La fundamental research, en cambio, encaja mejor en la universidad, porque puede prescindir por ahora de resultados de mercado y perseguir la coherencia interna de un sistema filosófico y una conciencia de problemas a largo plazo. El problema es que en la era de los LLM el umbral de recursos se ha elevado: incluso la fundamental research necesita, para verificar muchas hipótesis, una gran cantidad de cómputo y recursos de ingeniería que a la academia le resulta cada vez más difícil ofrecer.

\begin{importantbox}{El juicio de Zhao Chenyang (赵晨阳) sobre la applied research}
Lo más importante de la applied research es que debe tener una application. Si, al final, la dirección de un sistema debe demostrar su valor mediante la adopción de una gran empresa, usuarios reales y resultados comerciales, entonces formular una evaluation en el vacío dentro de la universidad puede resultar peor que ir directamente a la industria a recibir retroalimentación real.
\end{importantbox}

Aquí también hay una diferencia geográfica. Zhao Chenyang (赵晨阳) considera que a los estudiantes chinos que hacen el doctorado en Estados Unidos les cuesta más, porque las restricciones de visa y permiso de trabajo les dificultan mucho entrar, durante el doctorado, a las empresas más punteras y tener acceso a la tecnología más central. En cambio, muchas empresas de modelos en China están dispuestas a dar a sus becarios sin plaza formal permisos más altos; y en Hong Kong, Singapur o en un doctorado dentro de China tampoco existe necesariamente la misma necesidad urgente de irse. Por eso, la ola de abandono de estudios no es un consejo universal, sino el resultado de la superposición de ciertas regiones, ciertas identidades y ciertas ventanas tecnológicas.

El conductor añadió el matiz cultural de «rebeldía frente a lo ortodoxo»: en la concepción tradicional china, abandonar los estudios suele verse como algo poco honroso; mientras que en la narrativa del emprendimiento de Silicon Valley, abandonar los estudios puede entenderse más bien como un rasgo de outlier. Zhao Chenyang (赵晨阳) también admite que quienes están dispuestos a abandonar los estudios suelen ser menos apegados a las normas, y ese rasgo resulta atractivo tanto para emprender como para los inversionistas. Pero él no idealiza el abandono de estudios. Para mucha gente, quedarse en la universidad sigue siendo perfectamente razonable; lo que importa es si el problema en el que crees, tus condiciones de recursos y tu ventana de oportunidad todavía coinciden con el entorno universitario.

\begin{warningbox}{No hay que malinterpretar la ola de abandono de estudios como «el doctorado no sirve de nada»}
Lo que este episodio discute en realidad no es «si se debe o no abandonar los estudios», sino cuál es el entorno óptimo para sostener distintos tipos de research. La fundamental research todavía necesita a la universidad; la applied research de tipo sistema, si debe someterse de verdad a la adopción real y a la retroalimentación comercial, puede encajar mejor en la industria.
\end{warningbox}

Shang Yanyi (尚晏仪) cita la idea de Dai Yusen (戴雨森) de que emprender requiere cierta emoción y convicción. Si se calcula el desenlace final con total frialdad, quizá nunca se llegue a emprender, porque siempre se puede pensar que una gran empresa acabará por absorberte, que los eslabones de arriba y de abajo te presionarán, que el scope se irá reduciendo cada vez más. Emprender no surge de una planificación completamente racional; requiere cierta forma de confianza, impulso y disposición a asumir la incertidumbre. Esto también explica por qué, en cambio, los jóvenes pueden atreverse más a dar el salto: no están del todo fijados por un sistema de ciclos largos y aún conservan cierta preferencia por esa oportunidad del 1\%.

\subsection{Resumen de la sección}

La ola de abandono de estudios no es antiacademia, sino que, en la era de la IA, la relación entre la applied research y los recursos industriales, la retroalimentación real y la validación de mercado se vuelve más directa. La elección de Zhao Chenyang (赵晨阳) muestra que, si una dirección debe por naturaleza demostrarse mediante la adopción en la ingeniería y el valor comercial, entonces la industria puede ser más esencial que la universidad; pero la fundamental research, los problemas teóricos de largo plazo y los sistemas coherentes en sí mismos siguen teniendo valor en la universidad. Abandonar los estudios no es una plantilla, sino el resultado conjunto del timing, la identidad, los recursos y el tipo de problema.
\section{Research Taste en la era de la IA: lo que le queda al ser humano es el nivel de abstracción}

La parte intermedia de la conversación pasa de la deserción escolar a una pregunta más amplia: cuando la IA produce cada vez más artifacts, ¿cómo juzga el ser humano qué research es bueno? Zhao Chenyang (赵晨阳) y el conductor coinciden en que la applied research todavía puede reflejarse con revenue, valor para el usuario o benchmarks claros; la fundamental research es más difícil, porque se parece más a un arte y necesita personas con taste que la juzguen. En la era de la IA aumentan los artifacts, la capacidad de revisión humana es limitada y el peer review tradicional difícilmente resiste esa presión.

Aquí surge un juicio clave: el buen gusto científico se manifiesta en el nivel de abstracción con que se plantea un problema. Detrás de un problema concreto puede haber una estructura más fundamental, por ejemplo el consenso, la consistencia, la generalización o el aprendizaje a largo plazo dentro de la intelligence. La investigación que de verdad no es fácil de sustituir pronto por la IA no suele ser construir en seis meses un sistema soon-to-be-obsolete, sino lograr abstraer el problema de largo plazo que hay detrás de un testbed concreto.

\begin{importantbox}{El lugar del research taste}
Una vez que baja el costo de ejecución, el valor del ser humano no está en escribir más código o hacer más experimentos que la IA, sino en plantear abstracciones de problemas de más alto nivel, en juzgar qué problemas valen la pena, qué tradeoffs importan y qué rutas son solo optimizaciones locales.
\end{importantbox}

Zhao Chenyang comenta que le gusta escribir debug log y que mantiene un machine learning systems tutorial. Espera que en el futuro la IA pueda aprender de esos debug log el proceso de prueba y error. Este detalle es importante, porque los artículos académicos por lo general solo presentan la idea final ya pulida y rara vez registran los tradeoffs intermedios, los caminos fallidos y los juicios de depuración. Precisamente estos «juicios durante el proceso» quizá sean la razón por la que el taste es difícil de destilar y también difícil de sustituir con IA.

Esto se corresponde con el planteamiento de «la no linealidad de la investigación» de la entrevista de Xie Sanning (谢赛宁): cuando un artículo está escrito parece lineal, pero la idea real proviene de la exploración, el fracaso, los giros y la reinterpretación del problema. Si estos procesos no se registran, la IA solo puede aprender la forma del resultado y difícilmente puede aprender por qué el investigador abandonó cierta ruta, por qué conservó cierto diseño o por qué cierto fracaso en realidad aportaba mucha información.

\begin{knowledgebox}{Por qué importa el debug log}
\begin{itemize}
  \item Registra los caminos fallidos que un artículo no va a escribir.
  \item Expone los tradeoffs de diseño, en lugar de presentar solo la elección final.
  \item Convierte el «buen gusto», parcialmente, de una intuición indecible en material que se puede aprender.
  \item Puede llegar a ser una forma de datos importante para que la IA aprenda en el futuro el proceso de investigación.
\end{itemize}
\end{knowledgebox}

Este fragmento también señala una paradoja: una vez que la IA abarata la ejecución, un buen design se vuelve más importante. El costo de gobernar la escritura de código es mucho menor que el de contratar personas para ejecutar, siempre que el design sea suficientemente claro. La responsabilidad del ser humano ya no es escribir cada línea de código con sus propias manos, sino decidir la arquitectura, los límites, las prioridades, el manejo de fallos y los objetivos de largo plazo.

\subsection{Resumen de la sección}

En la era de la IA, el núcleo del research taste no es «escribir más», sino «abstraer más alto, juzgar con más precisión y registrar el proceso de manera más completa». La applied research puede apoyarse en la retroalimentación del mercado, mientras que la fundamental research depende más del taste y de una revisión de largo plazo. El debug log, el registro de tradeoffs y los caminos fallidos pueden llegar a ser material importante para que la IA futura aprenda el buen gusto en la investigación.
\section{AI Coding: después de que el código se abarata, la responsabilidad se encarece}

AI Coding es la línea más realista de este episodio. El anfitrión menciona que herramientas como Cloud Code hacen que la producción de código supere por mucho la velocidad humana, lo cual repercutirá en los reviewers de GitHub, los maintainers de código abierto y el code review interno de las empresas. Zhao Chenyang (赵晨阳) ya enfrenta dentro de SGLang la presión del código generado por IA, así que su respuesta es muy concreta.

Él no se opone a que la IA escriba código; al contrario, lo recibe con gusto. Lo esencial es esto: quien envía un PR debe poder explicar el design feature y estar dispuesto a responder por el resultado. Si alguien entrega un PR y ni siquiera puede aclarar los detalles ni la motivación del diseño, ese PR puede cerrarse sin ninguna contemplación. La IA no es un agente que asuma la responsabilidad por ti; si la usas para escribir código, cuando algo falle sigues siendo tú quien responde.

\begin{importantbox}{La regla central de los proyectos de código abierto frente a los PR de IA}
Se recibe con gusto que la IA escriba code, pero no que se haga sin pensar. El contribuyente debe explicar el diseño, entender los detalles y estar dispuesto a responder. Lo que la IA reduce es el costo de producir código, no la responsabilidad de mantener la calidad del proyecto.
\end{importantbox}

Este juicio tiene un sentido muy práctico. También se menciona el ejemplo de Amazon: se puede usar IA para escribir código, pero si algo sale mal hay que responder por ello. El mecanismo de responsabilidad termina por frenar el impulso de «generar 50K de código al día», porque cuanta más IA genera, más aumentan tanto la probabilidad de errores como el costo del review. Si un contribuyente solo busca acumular PR y registros de contribución, sin querer asumir la responsabilidad de mantenimiento a largo plazo del código, el proyecto de código abierto se desgastará hasta quebrarse.

El anfitrión plantea una observación muy importante: Git administra el resultado final del código, no el proceso de pensamiento. Vemos el diff final, pero no vemos cómo se formó ese diseño, cómo razonó el modelo, cómo interactuó el usuario con el coding agent, ni qué opciones se descartaron. Esto se conecta con el problema del debug log mencionado antes: lo verdaderamente escaso en la era de AI Coding no es solo el código en sí, sino el proceso de diseño, el nivel de pensamiento y la cadena de responsabilidad.

\begin{warningbox}{Git registra el código, no el juicio}
El gobierno del código en la era de la IA no puede basarse solo en el tamaño del diff, el número de PR o la cantidad de código. Dos PR con la misma cantidad de código pueden tener detrás un pensamiento, un riesgo y una responsabilidad de mantenimiento completamente distintos. Lo que en verdad habrá que registrar en el futuro es el decision flow, el proceso de review y las razones del diseño.
\end{warningbox}

Aquí se introduce la analogía entre el context graph y el SaaS. Lo verdaderamente valioso de un SaaS no es solo el materialized data, sino también cómo entraron esos datos al sistema, el proceso de negocio detrás y las offline discussion. Con Coding ocurre algo similar: el código final es solo el resultado; el flujo de diseño, el proceso de interacción y el contexto detrás son el activo más profundo. Quien logre capturar estos procesos podría formar, en la era de AI Coding, un nuevo ciclo cerrado de datos y capacidades.

\subsection{Resumen de la sección}

AI Coding abarata la producción de código, pero vuelve más importantes la responsabilidad, el diseño y el review. Los proyectos de código abierto no pueden resolver el problema prohibiendo el código de IA; deben establecer límites de responsabilidad: el contribuyente debe entender, explicar y mantener su propio PR. Los datos clave de la colaboración en código en el futuro tal vez no sean solo el diff final, sino el decision flow, el context graph y las razones del diseño.
\section{AI es un amplificador: de la retroalimentación de aprendizaje a la descarga cognitiva}

Zhao Chenyang (赵晨阳) dice que la IA es un amplificador: amplifica las demandas subjetivas y la iniciativa subjetiva de las personas. Este juicio tiene dos caras. Para quien aprende de manera activa, la IA es una bendición enorme. Por ejemplo, cuando en la preparatoria uno quería entender la regla de L'Hôpital, la enseñanza tradicional a veces solo decía «está fuera del temario, no se puede usar», sin explicar en serio cuándo se puede usar y por qué; la IA puede responder rápido a la curiosidad, dar retroalimentación inmediata y ayudar a explorar a fondo.

Pero Shang Yanyi (尚晏仪) advierte que hay que tener cuidado con la descarga cognitiva. Si una persona todavía no tiene claro qué quiere aprender o hacer, y le pasa el problema a la IA para que esta le entregue soluciones en varias rondas, a la larga perderá la capacidad de pensar de manera activa. La retroalimentación de la IA es muy oportuna, e incluso resulta complaciente con quien la usa, y la retroalimentación en imágenes y video es todavía más estimulante, por lo que es muy fácil establecer rápido un circuito de dopamina. Ella comenta que, al jugar con la generación de video por IA, sintió que este circuito de retroalimentación se establecía muy rápido, y por eso le pareció aterrador.

\begin{warningbox}{La IA como amplificador no amplifica automáticamente cosas buenas}
La IA amplifica la iniciativa subjetiva, pero también amplifica la pereza, la impulsividad y la adicción a la retroalimentación instantánea. Para quien ya tiene un objetivo claro, es un acelerador de aprendizaje; para quien no tiene objetivo ni capacidad de juicio, puede convertirse en una herramienta de descarga cognitiva.
\end{warningbox}

La discusión sobre el retiro de Sora y el modelo de negocio de la generación de video continúa con esta idea. Zhao Chenyang considera que la tecnología de Sora no es mala; el problema puede estar en que la inferencia de video es demasiado cara, y OpenAI no tiene un ecosistema de video propio que integre el contenido generado en el ciclo de revenue. Plataformas como Kuaishou o ByteDance, que ya cuentan con un ecosistema de video, pueden incorporar el video generado a su propio ciclo cerrado de distribución y monetización; si OpenAI hace un producto de video independiente, en cambio, enfrentará un descalce entre costo e ingreso.

Este pasaje en realidad extiende la idea de que «la IA es un amplificador» del aprendizaje individual al ecosistema comercial. El video de IA amplifica el deseo de entretenimiento de los usuarios y también la capacidad de distribución de las plataformas; pero sin un ciclo cerrado de ecosistema, que la tecnología sea fuerte no garantiza que sea buen negocio. La misma tecnología, en estructuras organizacionales distintas, produce resultados comerciales completamente diferentes.

\begin{knowledgebox}{Las tres capas de la lógica del amplificador}
\begin{tabularx}{\textwidth}{>{\raggedright\arraybackslash}p{0.2\textwidth} Y}
\toprule
Nivel & Qué amplifica \\
\midrule
Aprendizaje individual & Amplifica la curiosidad, la velocidad de retroalimentación y la capacidad de autoentrenamiento; también puede amplificar la descarga cognitiva. \\
Producción de código & Amplifica la velocidad de ejecución, también amplifica los PR de baja calidad y la presión de review. \\
Ecosistema de plataforma & Amplifica la capacidad de generación y distribución de contenido, pero solo si el costo puede integrarse a un ciclo comercial cerrado. \\
\bottomrule
\end{tabularx}
\end{knowledgebox}

\subsection{Resumen de la sección}

La IA no es simplemente una «herramienta más inteligente», sino un amplificador. Amplifica la exploración de quien aprende de manera activa, y también la dependencia de quien no piensa; amplifica la producción de código, y también la brecha de responsabilidad; amplifica la generación de video, y también el costo de las plataformas y el estímulo del contenido. Lo importante no es si se puede usar la IA o no, sino si quien la usa tiene un objetivo, un criterio estético y límites.
\section{El Agent no sabe decir que no: el límite de decisión humano}

En la segunda mitad de la conversación aparece un juicio muy importante: hoy es muy difícil que un Agent ayude a una persona a decidir «qué no hacer». Shang Yanyi (尚晏仪) pone como ejemplo el emprendimiento de producto: el objetivo es construir una profitable company, una meta relativamente cuantificable, pero las rutas de acción son infinitas y el entorno cambia sin cesar. Cada día se pueden hacer muchas cosas, y la competencia, la tecnología y las necesidades de los usuarios cambian constantemente. El Agent puede ayudar a ejecutar, pero es muy difícil que ayude a establecer prioridades, y en particular a decidir qué necesidades no atender.

Zhao Chenyang (赵晨阳) resumió el problema en una frase: la IA no say no, porque su reward consiste en satisfacer las necesidades humanas. Pero las necesidades de las personas no siempre son razonables. No toda expansión de negocio genera beneficios para la empresa; muchas necesidades, vistas a largo plazo, tienen un rendimiento negativo. En SGLang, la gran mayoría de las peticiones se rechazan sin miramientos, porque quienes mantienen el proyecto deben juzgar qué es correcto a largo plazo.

\begin{importantbox}{Decir que no es el núcleo de la capacidad de decisión humana}
La IA es hábil para satisfacer necesidades ya expresadas, pero el producto y la ingeniería a menudo requieren rechazar necesidades. Decidir qué hacer es importante, y decidir qué no hacer lo es igualmente. Los límites, las prioridades y el juicio de largo plazo son, por ahora, las responsabilidades humanas más difíciles de sustituir por un Agent.
\end{importantbox}

Esto se enlaza con el problema de responsabilidad de la IA Coding visto antes. Si una persona toma una decisión equivocada, la IA la ejecutará más rápido. Cuanto más barata es la ejecución, más rápido se amplifica también un rumbo equivocado. Por eso el criterio, el control de límites y la capacidad de decisión de las personas resultan aún más importantes. La era de la IA no significa que el ser humano se retire, sino que se repliega de la capa de ejecución hacia la capa de diseño, la capa de decisión y la capa de responsabilidad.

El moderador también conectó esta cuestión con el context graph y el decision flow: si en el futuro se quiere entrenar a un Agent para que tenga mejor juicio de producto, es necesario conservar los datos de cómo las personas toman decisiones. Pero hoy muchas de las discusiones clave ocurren fuera de línea, en reuniones, en mensajería instantánea y en el contexto de la organización, y es muy difícil capturarlas por completo. El código final o el producto final son solo el resultado; el verdadero «por qué se hizo así y por qué no de otra manera» a menudo se pierde.

\subsection{Resumen de la sección}

Una de las mayores limitaciones actuales del Agent no es no saber ejecutar, sino no saber negarse. Le resulta muy difícil comprender las prioridades bajo un objetivo de largo plazo, las restricciones de recursos, el costo de oportunidad y los límites estratégicos. La IA hace la ejecución más rápida, por lo que también los errores de decisión humanos se amplifican con mayor rapidez. En el futuro valdrá la pena registrar y entrenar no solo el resultado del código, sino también el decision flow y las razones para say no.

\section{Comunidad de código abierto: cómo registrar contribuciones y mantener el orden en la era de la IA}

El último tramo regresa a la comunidad de código abierto. Como proyecto de código abierto que recibe amplia atención, SGLang no enfrenta el problema de «si hay o no quién contribuya», sino que distintas etapas de la comunidad requieren distintas formas de gobernanza. En las primeras etapas quizá se acepten toda clase de solicitudes de funciones; una vez que el proyecto madura, alguien debe tomar la decisión final sobre qué se admite y qué no. De lo contrario, el proyecto termina arrastrado por demandas sin límite.

El código generado por IA agudiza este problema. El criterio de Zhao Chenyang (赵晨阳) es claro: se acepta que la IA escriba código, pero quien contribuye debe pensar, debe poder explicar los detalles del diseño y debe estar dispuesto a hacerse responsable. Acumular PRs, acumular contribuciones, entregar veinte PRs en un día sin estar dispuesto a entender ni a dar mantenimiento, daña al proyecto de código abierto. Quien contribuye de manera verdaderamente responsable, incluso usando IA, la trata como una herramienta de ejecución, no como una forma de externalizar la responsabilidad.

El anfitrión amplía la pregunta hacia la comunidad misma: cómo mantener una comunidad de código abierto, cómo reducir la asimetría de información, y cómo lograr que estudiantes e ingenieros jóvenes no solo reciban información, sino que también participen juntos en proyectos. La función básica de la comunidad es la transmisión de información; la función avanzada es la organización de la colaboración. Los meetups, los chats grupales, los proyectos de código abierto y los proyectos de evaluación pueden servir como vehículos para hacer las cosas en conjunto.

\begin{knowledgebox}{Tres niveles de función de una comunidad de código abierto}
\begin{tabularx}{\textwidth}{>{\raggedright\arraybackslash}p{0.2\textwidth} Y}
\toprule
Nivel & Función \\
\midrule
Transmisión de información & Permite que los nuevos integrantes conozcan más rápido lo que ocurre en el área, reduciendo la asimetría de información. \\
Organización de la colaboración & Permite que quienes tienen capacidad trabajen juntos en torno a proyectos, evaluaciones, herramientas y actividades. \\
Registro de contribuciones & Documenta quién hizo qué, la calidad de la contribución y cómo se hacen visibles las contribuciones que no son código. \\
\bottomrule
\end{tabularx}
\end{knowledgebox}

Hacia el final surge una pregunta muy interesante: cómo registrar la contribución de cada persona. La contribución en código es relativamente fácil de cuantificar, pero tampoco refleja del todo su naturaleza: la misma cantidad de código puede generarla distintos modelos, con una calidad y una cantidad de razonamiento del todo diferentes; las contribuciones que no son código, como organizar meetups, escribir documentación, hacer evaluaciones, difundir información o ayudar a los nuevos integrantes, son más difíciles de registrar. El anfitrión menciona que le gustaría construir una pequeña herramienta para consolidar y automatizar el registro de contribuciones. Esto coincide, de hecho, con el hilo conductor de toda la conversación: lo más difícil de registrar en la era de la IA no es el producto entregado, sino el proceso, el juicio y las relaciones.

\begin{warningbox}{La cantidad de contribuciones no es la calidad de las contribuciones}
La IA facilita fabricar cantidad de contribuciones, por lo que la comunidad no puede limitarse a mirar el número de PRs, las líneas de código o la frecuencia de los envíos. Lo que en verdad hay que medir es la comprensión del diseño, el mantenimiento a largo plazo, el juicio sobre los problemas, la calidad de la colaboración y el aporte al flujo de información de la comunidad.
\end{warningbox}

\subsection{Resumen de la sección}

Los proyectos de código abierto necesitan, en la era de la IA, una gobernanza más sólida, no fusiones más laxas. Puede aceptarse el código generado por IA, pero la responsabilidad no puede externalizarse. El valor de una comunidad no está solo en el repositorio de código, sino también en la transmisión de información, la organización de la colaboración, el registro de contribuciones y la formación de nuevos integrantes. La infraestructura de código abierto del futuro quizá se rediseñe en torno a «cómo registrar el proceso y la responsabilidad».
\section{Síntesis: la IA reduce el costo de ejecución, pero eleva el costo del juicio}

Toda la conversación puede resumirse en una frase: la IA está reduciendo rápidamente el costo de ejecución, pero eso no ha hecho desaparecer las responsabilidades humanas; al contrario, ha encarecido el juicio, la responsabilidad, el sentido estético y la capacidad de organización. La experiencia de Zhao Chenyang (赵晨阳) ofrece un buen corte transversal: pasó de la escuela a SGLang y luego a emprender con un motor de inferencia, enfrentando no solo oportunidades técnicas, sino también los límites de responsabilidad de la applied research, la adopción en la ingeniería, la gobernanza del código y la comunidad de código abierto.

\begin{longtable}{>{\raggedright\arraybackslash}p{0.24\textwidth} >{\raggedright\arraybackslash}p{0.68\textwidth}}
\toprule
Cuestión & Juicio central en la conversación \\
\midrule
\endfirsthead
\toprule
Cuestión & Juicio central en la conversación \\
\midrule
\endhead
Motor de inferencia & Se ubica entre la aplicación, el modelo y el hardware; es la capa intermedia que hace posible la operación de la IA a escala. \\
Ola de abandono escolar & En esencia no es un rechazo a la escuela, sino que la applied research necesita más aplicaciones reales, recursos y retroalimentación del mercado. \\
Applied vs Fundamental & La applied research necesita una application; la fundamental research necesita un sistema de preguntas autoconsistente. \\
Research taste & En la era de la IA, los seres humanos necesitan más abstracción de alto nivel, selección de problemas y juicio sobre tradeoffs. \\
Debug log & Los artículos ocultan el proceso de prueba y error; el debug log puede conservar un criterio de investigación que a la IA le cuesta aprender. \\
AI Coding & Una vez que producir código es barato, el diseño, el review, la responsabilidad y la capacidad de decir que no importan más. \\
Delegación cognitiva & La IA es un amplificador: puede amplificar la curiosidad, pero también la pereza y la adicción a la retroalimentación inmediata. \\
Decisiones del agente & El agente no dice que no por sí mismo; los seres humanos siguen teniendo que decidir qué no hacer. \\
Decision flow & El código final y el producto final son solo el resultado; los datos de juicio generados en el proceso son más escasos. \\
Comunidad de código abierto & En la era en que la IA genera código, el código abierto necesita registrar contribuciones, filtrar responsabilidades y mantener el orden del proyecto. \\
\bottomrule
\end{longtable}

Si el problema central de la generación anterior de la ingeniería de software era «cómo organizar el código que escriben los seres humanos», el problema de la era del AI Coding se convierte en «cómo reorganizar la producción amplificada por la IA, el juicio humano, los límites de responsabilidad y la colaboración comunitaria». Ahí radica el valor de esta conversación con Zhao Chenyang: no describe a la IA como una fuerza que simplemente sustituye a los seres humanos, sino que muestra cómo, después de que ocurre esa sustitución, emerge aquello que resulta más difícil de sustituir.
\section*{Apéndice: índice de juicios clave}

\begin{longtable}{>{\raggedright\arraybackslash}p{0.34\textwidth} >{\raggedright\arraybackslash}p{0.58\textwidth}}
\toprule
Juicio clave & Conciencia del problema correspondiente \\
\midrule
\endfirsthead
\toprule
Juicio clave & Conciencia del problema correspondiente \\
\midrule
\endhead
\textit{Lo más importante de la applied research es que tenga application.} & Esta es la razón central para dejar los estudios y elegir la industria: si la investigación se demuestra en definitiva mediante adoption real, el entorno industrial podría ser más esencial. \\
\textit{El círculo académico muchas veces sostiene un martillo buscando un clavo.} & Esta crítica apunta a la inercia de hacer formulate evaluation en el vacío, y también explica por qué gran parte del trabajo de sistemas necesita usuarios reales y retroalimentación real de ingeniería. \\
\textit{El buen gusto en la investigación científica se refleja en el nivel de abstracción del problema.} & La IA puede ejecutar tareas concretas más rápido; lo más escaso en los humanos es descubrir problemas de largo plazo y estructuras abstractas. \\
\textit{Me gusta escribir debug log.} & El debug log registra el ensayo y error, el tradeoff y el proceso de juicio que un artículo nunca escribiría; podría ser un material importante para que la IA del futuro aprenda el gusto por la investigación. \\
\textit{La IA es un amplificador.} & La IA amplifica la capacidad de actuar del ser humano, y también amplifica la descarga cognitiva, la adicción a la retroalimentación rápida y la producción de baja calidad. \\
\textit{La IA no say no.} & El Agent tiende a satisfacer las solicitudes, pero el producto, la ingeniería y el emprendimiento a menudo necesitan rechazar solicitudes. El juicio de límites sigue siendo responsabilidad humana. \\
\textit{Es bienvenido que la IA escriba code, pero tú debes poder explicarlo y responsabilizarte.} & Este es el límite mínimo de gobernanza que enfrentan los proyectos de código abierto ante el código generado por IA: la responsabilidad no se puede externalizar al modelo. \\
\textit{Git administra el código final, no el proceso de pensamiento.} & En la era de la AI Coding, el dato realmente escaso podría ser el decision flow, las razones de diseño y el proceso de interacción. \\
\textit{La comunidad no es solo transmisión de información, también es hacer cosas juntos.} & El siguiente paso de las comunidades de código abierto es pasar de reducir la asimetría de información a organizar proyectos, registrar contribuciones y acumular colaboración. \\
\bottomrule
\end{longtable}

\end{document}
